% ============================================================================
% Fichier  : partie1/P01_C02_grandes_defaillances.tex
% Titre    : Grandes Défaillances et Leçons de Gouvernance
% Auteur   : MINKA MI NGUIDJOÏ Thierry Emmanuel
% Version  : 1.0 -- octobre 2026
% ============================================================================
% Chapitre-cas : une section par affaire (brief, déroulé, révélations,
% synthèse), puis lecture croisée et grille d'erreurs.
% Sources : uniquement des documents publics de première main (décisions
% des régulateurs, analyses de cause publiées par les entreprises).
% ============================================================================

\chapter{Grandes Défaillances et Leçons de Gouvernance}
\label{chap:defaillances}
\chaptermark{Grandes défaillances}

\epigraph{%
    « Chaque grande panne informatique raconte deux histoires~: celle d'un
    composant qui a cédé, et celle d'une décision que personne n'a
    prise. »%
}{Principe d'analyse des défaillances}

\niveauChapitre{Fondamental}{%
    Maîtrise du Chapitre~\ref{chap:gouv-entreprise} (théorie de l'agence,
    gouvernance et management, organes de gouvernance).%
}

\begin{boxobjectifs}
\begin{itemize}
    \item Analyser des défaillances documentées à partir de sources
          publiques.
    \item Isoler, pour chaque cas, la défaillance de gouvernance derrière
          l'incident technique.
    \item Distinguer constatation et analyse dans un rapport de défaillance.
    \item Construire une grille d'erreurs de gouvernance réutilisable.
\end{itemize}
\end{boxobjectifs}

\bigskip

% ============================================================================
\section*{Mode d'emploi du chapitre}
% ============================================================================

Le prologue reposait sur un cas fictif, construit pour l'enseignement. Ce
chapitre passe aux faits. Il étudie trois défaillances réelles, choisies
parce qu'elles sont \textbf{documentées par des sources de première main}~:
la décision d'un régulateur des marchés, celle de deux régulateurs
bancaires, et l'analyse de cause publiée par un éditeur de logiciels. Elles
couvrent trois situations que tout DSI rencontrera~: un déploiement de
code, une migration de plateforme, et la mise à jour d'un logiciel fourni
par un tiers.

Chaque affaire suit le même plan~: le \textbf{brief} (qui, quand, quel
enjeu), le \textbf{déroulé} (la chronologie), les \textbf{révélations}
(ce que les enquêtes ont établi) et la \textbf{synthèse} (la cause de
gouvernance). La section~\ref{sec:c02-croisee} croise les trois cas, et
la section~\ref{sec:c02-grille} en tire une grille d'erreurs que vous
réutiliserez tout au long du cours.

\begin{boxalert}
\textbf{Règle de lecture.} Les faits rapportés dans ce chapitre sont ceux
qu'établissent les documents cités. Les analyses en termes de gouvernance
sont les nôtres~: elles sont signalées comme telles, avec leur niveau de
confiance. Les organisations citées ne sont pas jugées~; leurs défaillances
sont étudiées parce qu'elles ont eu le mérite d'être rendues publiques.
\end{boxalert}

\separateur

% ============================================================================
\section{Une Méthode d'Analyse des Défaillances}
\label{sec:c02-methode}
% ============================================================================

\subsection{Les sources}

Une analyse de défaillance vaut ce que valent ses sources. Les articles de
presse et les billets techniques sont utiles pour s'orienter, mais ils
simplifient et se recopient. Seules les sources de première main
permettent d'établir des constatations.

\begin{boxdef}
\termecle{Source de première main}\index{source de première main}
(\textit{primary source})~: document produit par une partie qui a eu
accès direct aux faits et qui engage sa responsabilité sur leur
exactitude~: décision d'un régulateur ou d'un tribunal, rapport d'enquête
indépendant commandé par le conseil, analyse de cause publiée par
l'entreprise elle-même.
\end{boxdef}

\begin{tableaucours}{L{3.6cm}L{4.2cm}L{4.4cm}}{%
    \tableauHeader{Type de source} &
    \tableauHeader{Force} &
    \tableauHeader{Limite}
}
Décision d'un régulateur & Faits contradictoirement établis, sanctions
motivées & Centrée sur les manquements à la réglementation \\
Rapport d'enquête indépendant & Examen large, accès aux pièces internes &
Commandé par l'entreprise, parfois publié en partie seulement \\
Analyse de cause de l'éditeur & Précision technique & Rédigée par la partie
en cause, prudente sur les responsabilités \\
Presse spécialisée & Rapidité, contexte & Chiffres approximatifs,
sources secondaires \\
\end{tableaucours}
\vspace{-0.8em}
\begin{center}{\small\itshape Tableau~2.1. Hiérarchie des sources d'une
analyse de défaillance}\end{center}

\subsection{La grille d'analyse}

Pour chaque affaire, sept questions structurent l'analyse. Elles reprennent
les réflexes du prologue (Chapitre~\ref{chap:prologue}).

\begin{enumerate}
    \item \textbf{Quel était l'enjeu~?} Activité, taille, dépendance au
          système.
    \item \textbf{Que s'est-il passé, et quand~?} Chronologie établie par
          les sources.
    \item \textbf{Quelle est la cause technique~?} Le composant, le code,
          la donnée.
    \item \textbf{Quelle décision a manqué~?} Qui décide, qui exécute, qui
          rend compte.
    \item \textbf{Quels contrôles auraient dû jouer~?} Avant, pendant et
          après l'incident.
    \item \textbf{Qui a demandé des comptes, et avec quelles
          conséquences~?} Régulateur, clients, marché, dirigeants.
    \item \textbf{Quelle leçon pour une organisation camerounaise~?}
\end{enumerate}

% ============================================================================
\section{Knight Capital, 2012 : 45 Minutes, 460 Millions de Dollars}
\label{sec:c02-knight}
% ============================================================================

\subsection{Brief}

Knight Capital Americas était, en 2012, l'un des principaux teneurs de
marché et intermédiaires sur les actions américaines. Son activité
reposait entièrement sur des systèmes automatisés d'acheminement des
ordres, dont le système SMARS. Les faits qui suivent sont ceux qu'établit
la décision de la Securities and Exchange Commission (\sigle{SEC}) du
16~octobre~2013~\cite{sec2013knight}.

\subsection{Déroulé}

Pour participer à un nouveau programme de la Bourse de New York, Knight
déploie un nouveau code sur les huit serveurs de SMARS dans les jours qui
précèdent le 1\textsuperscript{er}~août~2012. Ce nouveau code réutilise un
indicateur qui servait autrefois à une ancienne fonction, appelée « Power
Peg », abandonnée depuis des années mais jamais retirée du programme. Un
technicien déploie le code sur sept serveurs seulement~; le huitième garde
l'ancien code.

Le 1\textsuperscript{er}~août, à l'ouverture du marché, les ordres qui
portent le nouvel indicateur réveillent l'ancienne fonction sur le
huitième serveur. En 45~minutes, alors qu'il traitait 212~ordres de
clients, SMARS envoie des millions d'ordres sur le marché. Il en résulte
plus de 4~millions d'exécutions sur 154~valeurs, pour plus de
397~millions d'actions. Pour tenter de corriger, les équipes retirent le
nouveau code des sept serveurs où il avait été correctement installé, ce
qui étend le problème à l'ensemble des serveurs.

\subsection{Révélations}

La décision de la SEC établit notamment que~:
\begin{itemize}
    \item Knight ne disposait pas de procédures écrites de déploiement qui
          imposent une revue par un second technicien~;
    \item aucun contrôle ne vérifiait que le code était identique sur tous
          les serveurs~;
    \item avant l'ouverture du marché, le système avait émis des messages
          automatiques signalant une anomalie, que personne n'était chargé
          de traiter~;
    \item l'entreprise n'avait pas de procédure de réponse aux incidents
          adaptée à ce type de défaillance.
\end{itemize}
Knight a perdu plus de 460~millions de dollars sur les positions ainsi
créées. Elle a dû lever des fonds en urgence auprès d'investisseurs, puis
a été absorbée l'année suivante. La SEC lui a infligé une amende de
12~millions de dollars pour manquement à la règle sur l'accès au marché,
qui impose aux intermédiaires des contrôles de gestion des risques sur
leurs ordres automatisés.

\subsection{Synthèse}

\begin{boxscenario}{Knight Capital. Lecture de gouvernance}
\textbf{Cause technique} (établie). Un code obsolète, réactivé par un
indicateur réutilisé, sur un serveur où le déploiement était incomplet.

\medskip
\textbf{Cause de gouvernance} (analyse, confiance élevée). La gestion des
changements reposait sur le savoir-faire individuel, sans procédure
écrite, sans double contrôle et sans vérification après déploiement. La
dette technique (un code mort laissé en place) n'avait jamais été
arbitrée. Et le dispositif de surveillance produisait des alertes sans
qu'un rôle soit chargé de les traiter.

\medskip
\textbf{Les trois questions.} \emph{Qui décide} qu'un déploiement est
terminé~? Personne, formellement. \emph{Qui exécute}~? Un technicien seul.
\emph{Qui rend compte} des alertes avant l'ouverture~? Personne.
\end{boxscenario}

% ============================================================================
\section{TSB, 2018 : la Migration qui a Paralysé une Banque}
\label{sec:c02-tsb}
% ============================================================================

\subsection{Brief}

TSB Bank est une banque de détail britannique, rachetée en 2015 par le
groupe espagnol Sabadell. Elle utilisait jusqu'alors la plateforme
informatique de son ancien propriétaire. Le programme de migration vers
une nouvelle plateforme, fournie par une filiale informatique du groupe
Sabadell, devait lui donner son autonomie technique. Les faits qui
suivent sont ceux qu'établissent les décisions conjointes de la Financial
Conduct Authority (\sigle{FCA}) et de la Prudential Regulation Authority
(\sigle{PRA}) de décembre~2022~\cite{fca2022tsb}.

\subsection{Déroulé}

La migration des données des clients et de l'entreprise vers la nouvelle
plateforme a lieu lors d'un week-end d'avril~2018. Dès la remise en
service, toutes les agences et une part importante des 5,2~millions de
clients de la banque sont touchées~: accès impossibles aux services en
ligne et mobiles, opérations erronées, centres d'appel saturés. Le
rétablissement complet du service prend plusieurs mois. Le directeur
général quitte ses fonctions à l'automne~2018.

\subsection{Révélations}

Les régulateurs ont établi que TSB n'avait pas organisé ni contrôlé de
manière adéquate son programme de migration, et qu'elle n'avait pas géré
de manière adéquate les risques opérationnels nés de son externalisation
auprès d'un fournisseur tiers critique. La banque a versé 32,7~millions de
livres de réparation aux clients lésés. Les deux régulateurs lui ont
infligé une amende totale de 48,65~millions de livres (29,75~millions pour
la FCA, 18,9~millions pour la PRA), après une réduction de 30~\% accordée
pour son règlement amiable. En 2023, la PRA a en outre sanctionné
personnellement l'ancien directeur des systèmes d'information de la banque
pour ne pas avoir pris de mesures raisonnables de supervision du
prestataire~\cite{pra2023abarca}.

\begin{boxinfo}
\textbf{Un précédent pour tous les DSI.} La sanction personnelle de
l'ancien DSI de TSB montre que la redevabilité n'est pas une notion
abstraite. Dans le régime britannique de responsabilité des dirigeants du
secteur financier, le titulaire d'une fonction clé répond personnellement
des mesures raisonnables qu'il a prises, ou omis de prendre, dans son
périmètre. Externaliser une activité ne transfère pas la responsabilité de
la superviser.
\end{boxinfo}

\subsection{Synthèse}

\begin{boxscenario}{TSB. Lecture de gouvernance}
\textbf{Cause technique} (établie dans ses grandes lignes). Une plateforme
et un fournisseur qui n'étaient pas prêts à absorber la charge et les cas
réels de l'ensemble des clients au moment de la bascule.

\medskip
\textbf{Cause de gouvernance} (analyse, confiance élevée). Le programme n'a
pas été piloté avec un contrôle suffisant de son état de préparation, et
la banque s'est fiée à son fournisseur, entité de son propre groupe, sans
superviser elle-même les risques qu'il portait. La décision de basculer a
été prise sans que l'organe de décision dispose d'une évaluation
indépendante de cet état de préparation.

\medskip
\textbf{Les trois questions.} \emph{Qui décide} de la bascule~? La
direction, sur la foi des assurances du programme. \emph{Qui exécute}~?
Un fournisseur dont la banque ne maîtrisait pas les risques. \emph{Qui
rend compte}~? La banque, son directeur général et son DSI, devant les
régulateurs.
\end{boxscenario}

La ressemblance avec l'opération SANAGA n'est pas fortuite~: le cas fictif
du prologue reprend des mécanismes que les décisions des régulateurs
britanniques ont rendus publics. Ils ne sont propres ni à une banque ni à
un pays.

% ============================================================================
\section{CrowdStrike, 2024 : la Mise à Jour qui a Arrêté le Monde}
\label{sec:c02-crowdstrike}
% ============================================================================

\subsection{Brief}

CrowdStrike est un éditeur américain de solutions de cybersécurité. Son
logiciel Falcon protège les postes de travail et les serveurs de très
nombreuses organisations. Pour détecter les attaques, il s'exécute au cœur
du système d'exploitation Windows et reçoit plusieurs fois par jour des
mises à jour de contenu, qui décrivent de nouveaux comportements
malveillants. Les faits techniques qui suivent sont ceux qu'établit
l'analyse de cause publiée par l'éditeur le 6~août~2024~\cite{crowdstrike2024rca}.

\subsection{Déroulé}

Le 19~juillet~2024, CrowdStrike diffuse une mise à jour de contenu pour
son capteur Windows. Elle provoque l'arrêt brutal des machines qui la
reçoivent, puis leur redémarrage en boucle. L'éditeur retire la mise à
jour dans l'heure et demie qui suit, mais les machines déjà touchées ne
peuvent plus la recevoir~: il faut souvent intervenir manuellement, poste
par poste. Microsoft estime à environ 8,5~millions le nombre d'appareils
Windows touchés~\cite{microsoft2024crowdstrike}. Des compagnies aériennes,
des hôpitaux, des banques et des médias sont paralysés pendant plusieurs
jours.

\subsection{Révélations}

L'analyse de cause établit que le capteur attendait 20~champs de données,
alors que la mise à jour en fournissait 21. Ce décalage a provoqué une
lecture hors des limites de la mémoire, et l'arrêt du système. Le défaut
n'avait pas été détecté par les contrôles de validation du contenu.
L'éditeur annonce, parmi ses mesures correctives, davantage de tests, un
déploiement progressif de ces mises à jour, et la possibilité pour les
clients de choisir où et quand elles s'installent.

\subsection{Synthèse}

Ce cas se lit à deux niveaux de gouvernance~: celle de l'éditeur, et celle
de ses clients.

\begin{boxscenario}{CrowdStrike. Lecture de gouvernance}
\textbf{Chez l'éditeur} (analyse, confiance moyenne). Les mises à jour de
contenu n'étaient pas soumises à la même rigueur de gestion des
changements que les versions du logiciel, ni déployées progressivement.
Un type de changement fréquent et jugé peu risqué échappait aux contrôles
les plus stricts.

\medskip
\textbf{Chez les clients} (analyse, confiance élevée). Des milliers
d'organisations avaient accordé à un fournisseur unique le droit de
modifier, sans délai ni contrôle de leur part, un composant qui s'exécute
au cœur de tous leurs postes. Ce choix est rationnel pour la sécurité~:
une protection doit être à jour. Mais il crée une \termecle{dépendance
critique}\index{dépendance critique} qui n'avait, dans la plupart des cas,
jamais été évaluée ni arbitrée par les instances de gouvernance.
\end{boxscenario}

\begin{boxmetier}{RSSI}
\textbf{Évaluer une dépendance critique à un fournisseur de sécurité}
\begin{enumerate}
    \item Recenser les logiciels qui disposent de droits élevés sur
          l'ensemble du parc et reçoivent des mises à jour automatiques.
    \item Pour chacun, vérifier si l'organisation peut retarder ou
          échelonner les mises à jour (groupe pilote, puis déploiement
          général).
    \item Évaluer le scénario « le fournisseur rend tout le parc
          inutilisable »~: délai de restauration, procédure manuelle,
          personnel nécessaire.
    \item Présenter ce risque au comité des risques, avec une
          recommandation~: l'accepter, le réduire ou le transférer.
\end{enumerate}
Ce qui compte n'est pas de supprimer la dépendance, souvent inévitable,
mais qu'elle soit \textbf{connue et décidée} par l'instance compétente.
\end{boxmetier}

% ============================================================================
\section{Lecture Croisée des Trois Affaires}
\label{sec:c02-croisee}
% ============================================================================

\subsection{Comparaison}

\begin{tableaucours}{L{2.9cm}L{3.1cm}L{3.1cm}L{3.1cm}}{%
    \tableauHeader{Question} &
    \tableauHeader{Knight Capital} &
    \tableauHeader{TSB} &
    \tableauHeader{CrowdStrike}
}
Nature du changement & Déploiement de code & Migration de plateforme &
Mise à jour de contenu \\
Durée de l'incident & 45 minutes & Plusieurs mois & Plusieurs jours \\
Contrôle défaillant & Déploiement sans procédure ni double contrôle &
Préparation de la bascule, supervision du fournisseur & Validation et
déploiement progressif \\
Rôle d'un tiers & Aucun & Fournisseur critique du groupe & L'éditeur
lui-même est le tiers des clients \\
Sanction & Amende de la SEC, absorption & Amendes FCA et PRA, sanction du
DSI & Pertes et contentieux des clients \\
\end{tableaucours}
\vspace{-0.8em}
\begin{center}{\small\itshape Tableau~2.2. Comparaison des trois affaires}\end{center}

\subsection{Quatre constantes}

Malgré leurs différences, les trois affaires partagent quatre traits.

\begin{enumerate}
    \item \textbf{Un changement est à l'origine de chaque incident.} Aucune
          de ces pannes ne vient d'une usure ou d'une attaque. Toutes
          viennent d'une modification volontaire du système. La gestion
          des changements (Chapitre~\ref{chap:itil-changements}) est le
          premier contrôle de la gouvernance des SI.
    \item \textbf{Le risque était connu ou connaissable.} Code mort jamais
          retiré, préparation insuffisante, type de mise à jour moins
          contrôlé~: aucun de ces risques n'était imprévisible.
    \item \textbf{La responsabilité était diluée.} Dans chaque cas, il est
          difficile de nommer la personne ou l'instance qui avait
          formellement décidé que le risque était acceptable.
    \item \textbf{Le retour arrière était difficile ou impossible.} Knight a
          aggravé la situation en tentant de revenir en arrière, TSB a mis
          des mois à se rétablir, et les clients de CrowdStrike ont dû
          réparer leurs postes un par un.
\end{enumerate}

\begin{boiteRemarque}{Et en Afrique~?}
Les trois affaires sont américaines ou européennes. Ce n'est pas que
l'Afrique centrale ignore les grandes pannes~: les interruptions de
services bancaires, de paiement mobile ou de télécommunications y sont
fréquentes et parfois longues. C'est que leurs causes sont rarement
rendues publiques. Peu de régulateurs publient leurs décisions motivées,
et peu d'entreprises publient leurs analyses de cause.

Cette opacité est elle-même une question de gouvernance. Une défaillance
dont on ne publie pas l'analyse ne profite à personne d'autre qu'à
l'organisation qui l'a subie, et souvent pas même à elle. Le
Chapitre~\ref{chap:droit} examinera les obligations de déclaration des
incidents prévues par le droit camerounais, et le
Chapitre~\ref{chap:recherche-mesure} la place de cette question dans la
recherche.
\end{boiteRemarque}

% ============================================================================
\section{Une Grille d'Erreurs de Gouvernance}
\label{sec:c02-grille}
% ============================================================================

Les trois affaires réelles et l'opération SANAGA permettent de construire
une grille de dix erreurs. Chacune s'accompagne d'un signal d'alerte
observable et de la question qu'un administrateur, un auditeur ou un DSI
devrait poser. La dernière colonne renvoie à l'objectif COBIT~2019 qui
traite l'erreur~; ces objectifs seront étudiés au
Chapitre~\ref{chap:cobit-objectifs}.

\begin{tableaucours}{L{3.2cm}L{3.4cm}L{3.8cm}L{1.6cm}}{%
    \tableauHeader{Erreur} &
    \tableauHeader{Signal d'alerte} &
    \tableauHeader{Question à poser} &
    \tableauHeader{COBIT}
}
1. Décision sans instance & Date ou budget fixés par une seule personne &
Quelle instance a décidé, et où est le procès-verbal~? & EDM01 \\
2. Risque non arbitré & Alerte écrite restée sans réponse & Qui a accepté
ce risque, par écrit~? & EDM03 \\
3. Changement non autorisé & Mise en production sans comité & Où est
l'autorisation, datée avant la mise en production~? & BAI06 \\
4. Pas de retour arrière & Plan de retour absent ou non testé & Quand le
retour arrière a-t-il été testé, et avec quel résultat~? & BAI07 \\
5. Recette incomplète & Bascule avec des scénarios non testés & Quel
taux de couverture, et qui a accepté le reste~? & BAI07 \\
6. Fournisseur non supervisé & Rapports de suivi rédigés par le
fournisseur & Comment vérifions-nous nous-mêmes ce qu'il affirme~? &
APO10 \\
7. Contrat sans engagement & Pas de niveau de service ni de pénalité &
Que se passe-t-il s'il ne tient pas ses engagements~? & APO09 \\
8. Alertes sans propriétaire & Messages automatiques non traités & Qui est
chargé de lire cette alerte, et dans quel délai~? & DSS02 \\
9. Continuité non testée & Plan de continuité jamais exercé & Quand a eu
lieu le dernier exercice, et qu'a-t-il révélé~? & DSS04 \\
10. Audit sans suite & Recommandations « en cours » depuis des mois &
Qui suit les recommandations, devant quelle instance~? & MEA02 \\
\end{tableaucours}
\vspace{-0.8em}
\begin{center}{\small\itshape Tableau~2.3. Grille des dix erreurs de
gouvernance des SI}\end{center}

\begin{boxprocedure}
\textbf{Le dossier d'analyse d'une défaillance.} Pour qu'une analyse de
défaillance soit opposable à la direction, au conseil ou au régulateur,
elle doit contenir~:
\begin{itemize}
    \item la liste des sources utilisées, classées selon le
          Tableau~2.1~;
    \item une chronologie où chaque entrée renvoie à une pièce~;
    \item des constatations séparées des analyses, chaque analyse portant
          un niveau de confiance (faible, moyen, élevé)~;
    \item la cause technique et la cause de gouvernance, distinguées~;
    \item pour chaque erreur de la grille relevée, la pièce qui l'établit.
\end{itemize}
\end{boxprocedure}

\begin{boxmetier}{Auditeur SI}
\textbf{Utiliser la grille en mission}
\begin{enumerate}
    \item Avant la mission, choisir les erreurs de la grille les plus
          probables au regard de l'activité (une banque qui migre~:
          erreurs 1 à 5~; une entreprise très externalisée~: erreurs 6
          et~7).
    \item Pour chaque erreur retenue, préparer la question et la pièce à
          demander.
    \item En entretien, poser la question, puis demander la pièce~:
          une réponse orale n'est pas une preuve.
    \item Consigner chaque écart sous la forme constatation, critère,
          cause, conséquence, recommandation
          (Chapitre~\ref{chap:audit}).
\end{enumerate}
\end{boxmetier}

\begin{boxmetier}{Responsable des services IT}
\textbf{Trois contrôles à mettre en place dès demain}
\begin{enumerate}
    \item Aucun changement en production sans enregistrement préalable, et
          aucun changement majeur sans autorisation d'un comité.
    \item Une vérification après chaque déploiement~: la version installée
          est la même partout.
    \item Un propriétaire nommé pour chaque alerte automatique, avec un
          délai de traitement.
\end{enumerate}
Ces trois contrôles auraient, à eux seuls, très probablement évité
l'affaire Knight Capital.
\end{boxmetier}

% ============================================================================
\separateur
% ============================================================================

\begin{boxresume}
\textbf{Ce qu'il faut retenir de ce chapitre~:}
\begin{itemize}
    \item \textbf{Les sources de première main d'abord}~: décisions des
          régulateurs, rapports d'enquête, analyses de cause. La presse
          oriente, elle n'établit pas.
    \item \textbf{Knight Capital (2012)}~: un déploiement sans procédure
          écrite ni double contrôle, un code mort réactivé, des alertes
          sans propriétaire~; 460~millions de dollars perdus en
          45~minutes.
    \item \textbf{TSB (2018)}~: une migration mal pilotée, un fournisseur
          critique mal supervisé~; des amendes de 48,65~millions de livres
          et la sanction personnelle de l'ancien DSI.
    \item \textbf{CrowdStrike (2024)}~: une mise à jour de contenu
          insuffisamment validée et déployée en une fois~; une dépendance
          critique que la plupart des clients n'avaient jamais arbitrée.
    \item \textbf{Quatre constantes}~: un changement à l'origine, un risque
          connaissable, une responsabilité diluée, un retour arrière
          difficile.
    \item \textbf{La grille des dix erreurs} relie chaque erreur à un
          signal, une question et un objectif COBIT~: elle servira dans
          tout le cours.
\end{itemize}
\end{boxresume}

\bigskip

\begin{boxexo}[title={\ding{45}\quad Exercice~2.1. Constatation ou analyse~? \niveaubase}]
Pour chacune des phrases suivantes, indiquez s'il s'agit d'une constatation
ou d'une analyse. Pour les analyses, proposez un niveau de confiance et la
pièce qui permettrait de le renforcer.
\begin{enumerate}[(a)]
    \item « Le nouveau code n'a été installé que sur sept des huit
          serveurs. »
    \item « Les équipes de Knight manquaient de formation. »
    \item « TSB a versé 32,7~millions de livres de réparation à ses
          clients. »
    \item « Le conseil de TSB aurait dû refuser la date de bascule. »
    \item « La mise à jour fournissait 21~champs au lieu des 20 attendus. »
\end{enumerate}
\end{boxexo}

\begin{boxexo}[title={\ding{45}\quad Exercice~2.2. La grille appliquée à KRIBILOG \niveaumoyen}]
Chez KRIBILOG SA, le logiciel de suivi des conteneurs s'arrête trois fois
en deux mois. Vous disposez des faits suivants~: l'infogérant NETPORT
SERVICES applique les correctifs le soir, sans prévenir~; son rapport
mensuel ne mentionne aucun incident~; le contrat ne prévoit aucun délai
d'intervention~; le dernier test de continuité date de 2021.
\begin{enumerate}[(a)]
    \item Identifiez les erreurs de la grille (Tableau~2.3) présentes dans
          ce cas.
    \item Pour chacune, rédigez la question à poser et la pièce à
          demander.
    \item Classez-les par ordre de priorité, en justifiant.
\end{enumerate}
\end{boxexo}

\begin{boxexo}[title={\ding{45}\quad Exercice~2.3. Simulation d'examen \niveauexpert}]
\textit{Durée conseillée~: 1~heure.} Choisissez une interruption de service
informatique survenue en Afrique et documentée publiquement (communiqué de
l'organisation, décision d'un régulateur, rapport d'enquête).
\begin{enumerate}[(a)]
    \item Classez les sources disponibles selon le Tableau~2.1, et
          indiquez ce qu'elles permettent d'établir.
    \item Appliquez les sept questions de la grille d'analyse
          (section~\ref{sec:c02-methode}).
    \item Rédigez une synthèse d'une page, en séparant constatations et
          analyses.
    \item Si les sources ne permettent pas d'établir la cause de
          gouvernance, dites-le, et indiquez les pièces qu'il faudrait
          obtenir. Cette conclusion est une réponse valable.
\end{enumerate}
\end{boxexo}

\bigskip

\begin{boxinfo}
\textbf{Transition.} Les défaillances montrent ce que coûte une
gouvernance absente. Le Chapitre~\ref{chap:alignement} montre ce qu'une
gouvernance présente doit rechercher~: l'alignement du SI sur la
stratégie, la création de valeur, et l'arbitrage des décisions
d'investissement.
\end{boxinfo}
